\section{\texorpdfstring{\(RG\)-Lattices and Lifting of Projectives}{RG-Lattices and Lifting of Projectives}}

\begin{definition}[\texorpdfstring{\(RG\)-lattice}{RG-lattice}]
    Let \(R\) be a commutative ring and let \(G\) be a finite group.
    An \(RG\)-module \(L\) is called an \textbf{\(RG\)-lattice} if
    \(L\) is finitely generated as an \(RG\)-module and is free as an \(R\)-module.
\end{definition}

\begin{proposition}[Lifting projective covers from \texorpdfstring{\(kG\)}{kG} to \texorpdfstring{\(RG\)}{RG}]
    Let \(R\) be a complete discrete valuation ring with maximal ideal
    \[
        J(R)=(\pi),
    \]
    and let
    \[
        k=R/(\pi).
    \]
    Let \(G\) be a finite group. Then the following hold.

    \begin{enumerate}
        \item For every simple \(kG\)-module \(S\), there exists an indecomposable
        projective \(RG\)-module \(\widehat P_S\) such that
        \[
            \widehat P_S=RG\widehat e_S,
        \]
        where \(\widehat e_S\in RG\) is a primitive idempotent satisfying
        \[
            \widehat e_S S\neq 0.
        \]
        Moreover, \(\widehat P_S\) is a projective cover of \(S\) as an \(RG\)-module.

        \item If \(P_S\) denotes the projective cover of \(S\) as a \(kG\)-module, then
        \[
            \widehat P_S/\pi\widehat P_S\cong P_S.
        \]
        Moreover, the canonical epimorphism
        \[
            \widehat P_S\longrightarrow \widehat P_S/\pi\widehat P_S\cong P_S
        \]
        is a projective cover of \(P_S\) as an \(RG\)-module. In particular,
        \(S\) is the unique simple quotient of \(\widehat P_S\).

        Consequently, for simple \(kG\)-modules \(S\) and \(T\),
        \[
            \widehat P_S\cong \widehat P_T
            \quad\Longleftrightarrow\quad
            S\cong T.
        \]

        \item Every finitely generated \(RG\)-module has a projective cover.

        \item Every finitely generated indecomposable projective \(RG\)-module is
        isomorphic to \(\widehat P_S\) for some simple \(kG\)-module \(S\).

        \item Every finitely generated projective \(kG\)-module can be lifted to a
        projective \(RG\)-lattice. Such a lift is unique up to isomorphism.
    \end{enumerate}
\end{proposition}

\begin{proof}
    Let
    \[
        \overline{RG}=RG/\pi RG\cong kG.
    \]
    Since \(R\) is complete, idempotents lift modulo the ideal \(\pi RG\). Moreover,
    \[
        \pi RG\subseteq J(RG),
    \]
    so primitive idempotents lift to primitive idempotents.

    Let \(S\) be a simple \(kG\)-module. Choose a primitive idempotent
    \(e_S\in kG\) such that
    \[
        P_S=kG e_S
    \]
    is the projective cover of \(S\), and
    \[
        e_S S\neq 0.
    \]
    Lift \(e_S\) to a primitive idempotent
    \[
        \widehat e_S\in RG.
    \]
    Define
    \[
        \widehat P_S:=RG\widehat e_S.
    \]
    Since \(\widehat e_S\) is an idempotent, \(RG\widehat e_S\) is a direct summand
    of the regular module \(RG\), hence is projective. Since \(\widehat e_S\) is
    primitive, \(RG\widehat e_S\) is indecomposable.

    Reducing modulo \(\pi\), we get
    \[
        \widehat P_S/\pi\widehat P_S
        \cong
        (RG/\pi RG)(\widehat e_S+\pi RG)
        \cong
        kG e_S
        =
        P_S.
    \]
    Thus
    \[
        \widehat P_S/\pi\widehat P_S\cong P_S.
    \]

    Since
    \[
        \pi\widehat P_S\subseteq \operatorname{Rad}(\widehat P_S),
    \]
    the canonical epimorphism
    \[
        \widehat P_S\longrightarrow \widehat P_S/\pi\widehat P_S
    \]
    is essential. Hence
    \[
        \widehat P_S\longrightarrow P_S
    \]
    is a projective cover of \(P_S\) as an \(RG\)-module.

    Composing with the projective cover
    \[
        P_S\longrightarrow S
    \]
    in \(kG\)-modules, we obtain an essential epimorphism
    \[
        \widehat P_S\longrightarrow S.
    \]
    Therefore \(\widehat P_S\) is a projective cover of \(S\) as an \(RG\)-module.

    Moreover,
    \[
        \operatorname{Top}(\widehat P_S)
        =
        \widehat P_S/\operatorname{Rad}(\widehat P_S)
        \cong
        P_S/\operatorname{Rad}(P_S)
        \cong
        S.
    \]
    Hence \(S\) is the unique simple quotient of \(\widehat P_S\). It follows that
    for simple \(kG\)-modules \(S\) and \(T\),
    \[
        \widehat P_S\cong \widehat P_T
        \quad\Longleftrightarrow\quad
        \operatorname{Top}(\widehat P_S)\cong \operatorname{Top}(\widehat P_T)
        \quad\Longleftrightarrow\quad
        S\cong T.
    \]

    Now let \(U\) be a finitely generated \(RG\)-module. Since \(RG/J(RG)\) is a
    semisimple Artinian ring, the top
    \[
        \operatorname{Top}(U)=U/\operatorname{Rad}(U)
    \]
    is a semisimple \(RG\)-module. Hence
    \[
        \operatorname{Top}(U)
        \cong
        S_1\oplus\cdots\oplus S_m
    \]
    for some simple \(kG\)-modules \(S_i\). Then
    \[
        \widehat P_{S_1}\oplus\cdots\oplus \widehat P_{S_m}
        \longrightarrow
        S_1\oplus\cdots\oplus S_m
        \cong
        \operatorname{Top}(U)
    \]
    is a projective cover. Since the canonical map
    \[
        U\longrightarrow \operatorname{Top}(U)
    \]
    is an essential epimorphism, lifting along it gives an essential epimorphism from
    a projective module onto \(U\). Thus \(U\) has a projective cover.

    Let \(P\) be a finitely generated indecomposable projective \(RG\)-module. The
    canonical epimorphism
    \[
        P\longrightarrow \operatorname{Top}(P)
    \]
    is a projective cover. Since \(\operatorname{Top}(P)\) is semisimple, write
    \[
        \operatorname{Top}(P)\cong S_1\oplus\cdots\oplus S_m.
    \]
    Then the projective cover of \(\operatorname{Top}(P)\) is
    \[
        \widehat P_{S_1}\oplus\cdots\oplus \widehat P_{S_m}.
    \]
    By uniqueness of projective covers,
    \[
        P\cong \widehat P_{S_1}\oplus\cdots\oplus \widehat P_{S_m}.
    \]
    Since \(P\) is indecomposable, we must have \(m=1\). Hence
    \[
        P\cong \widehat P_S
    \]
    for some simple \(kG\)-module \(S\).

    Finally, let \(Q\) be a finitely generated projective \(kG\)-module. Since
    \(kG\) is Artinian, \(Q\) decomposes as a finite direct sum of indecomposable
    projective \(kG\)-modules:
    \[
        Q\cong \bigoplus_S P_S^{\oplus m_S}.
    \]
    Define
    \[
        \widehat Q:=\bigoplus_S \widehat P_S^{\oplus m_S}.
    \]
    Then \(\widehat Q\) is a projective \(RG\)-lattice and
    \[
        \widehat Q/\pi\widehat Q
        \cong
        \bigoplus_S P_S^{\oplus m_S}
        \cong
        Q.
    \]
    Thus \(Q\) lifts to a projective \(RG\)-lattice. The uniqueness follows from the
    isomorphism criterion below.
\end{proof}

\begin{corollary}
    Let \(L\) be an \(RG\)-lattice. Then \(L\) is projective as an \(RG\)-module if
    and only if
    \[
        L/\pi L
    \]
    is projective as a \(kG\)-module.
\end{corollary}

\begin{proof}
    If \(L\) is projective over \(RG\), then \(L\) is a direct summand of a finite free
    \(RG\)-module. Reducing modulo \(\pi\), we see that \(L/\pi L\) is a direct
    summand of a finite free \(kG\)-module. Hence \(L/\pi L\) is projective over
    \(kG\).

    Conversely, assume that \(L/\pi L\) is projective over \(kG\). Let
    \[
        P\longrightarrow L
    \]
    be a projective cover of \(L\) as an \(RG\)-module. Reducing modulo \(\pi\), we
    obtain an essential epimorphism
    \[
        P/\pi P\longrightarrow L/\pi L.
    \]
    Since \(L/\pi L\) is projective, this essential epimorphism must be an isomorphism.

    Let \(K\) be the kernel of \(P\to L\). Since \(P\) and \(L\) are free over the DVR
    \(R\), reduction modulo \(\pi\) gives
    \[
        K/\pi K=0.
    \]
    By Nakayama's lemma, \(K=0\). Hence
    \[
        P\cong L,
    \]
    and therefore \(L\) is projective.
\end{proof}

\begin{corollary}
    Let \(L_1\) and \(L_2\) be finitely generated projective \(RG\)-modules. Then
    \[
        L_1\cong L_2
        \quad\Longleftrightarrow\quad
        L_1/\pi L_1\cong L_2/\pi L_2
    \]
    as \(kG\)-modules.
\end{corollary}

\begin{proof}
    The forward implication is immediate by reduction modulo \(\pi\).

    Conversely, suppose
    \[
        L_1/\pi L_1\cong L_2/\pi L_2
    \]
    as \(kG\)-modules. Since \(L_i\) is projective over \(RG\), the natural map
    \[
        L_i\longrightarrow L_i/\pi L_i
    \]
    is a projective cover. Therefore \(L_1\) and \(L_2\) are projective covers of
    isomorphic \(RG\)-modules. By uniqueness of projective covers,
    \[
        L_1\cong L_2.
    \]
\end{proof}

